% Motor: XeLaTeX. Sesión 27 - material docente UTP.
% Estructura docente de la Sesión 26.
\documentclass[aspectratio=169,10pt,c]{beamer}
% Reutiliza sin cambios la identidad visual UTP validada en la Sesion 08.
% Reutiliza sin cambios la identidad visual UTP validada en la Sesion 05.
% La segunda ruta permite compartir los assets comunes sin duplicar binarios.
% !TEX root = slides.tex
% Preambulo comun: estilo UTP, colores, fuentes, plantillas de beamer y
% comandos de apoyo didactico. Se carga desde slides.tex tras \documentclass.

% La opción "c" de beamer reparte el espacio libre 1:1.5 (arriba:abajo), por lo
% que el contenido queda por encima del centro. Igualamos ambos rellenos.
\makeatletter
\define@key{beamerframe}{c}[true]{%
  \beamer@frametopskip=0pt plus 1fill\relax%
  \beamer@framebottomskip=0pt plus 1fill\relax%
  \beamer@frametopskipautobreak=0pt plus .5\paperheight\relax%
  \beamer@framebottomskipautobreak=0pt plus .5\paperheight\relax%
}
\makeatother

\usepackage{graphicx}
\usepackage{tikz}
\usepackage{xcolor}
\usepackage{iftex}

\ifPDFTeX
  \usepackage[utf8]{inputenc}
  \usepackage[T1]{fontenc}
  \usepackage{lmodern}          % tipografia escalable para el respaldo con pdflatex
  \renewcommand{\familydefault}{\sfdefault}
\else
  \usepackage{fontspec}
  \IfFontExistsTF{Arial}
    {\setsansfont{Arial}[BoldFont={Arial Bold}]}
    {\setsansfont{Helvetica}[BoldFont={Helvetica Bold}]}
\fi

\usepackage[spanish,es-nodecimaldot]{babel}
\usepackage{microtype}
\usepackage{booktabs}
\usepackage{array}
\usetikzlibrary{arrows.meta,positioning,fit,calc,shapes.geometric}

\graphicspath{{assets/}}

% ---------- Identidad visual UTP ----------
\definecolor{utpRojo}{HTML}{C42B35}
\definecolor{utpNegro}{HTML}{111111}
\definecolor{utpGris}{HTML}{5F5F5F}
\definecolor{utpGrisMedio}{HTML}{9A9A9A}
\definecolor{utpGrisClaro}{HTML}{F1F1F1}

\usetheme{default}
\setbeamercolor{normal text}{fg=utpNegro,bg=white}
\setbeamercolor{structure}{fg=utpRojo}
\setbeamercolor{frametitle}{fg=utpNegro,bg=white}
\setbeamercolor{title}{fg=utpNegro}
\setbeamercolor{subtitle}{fg=utpNegro}
\setbeamercolor{author}{fg=utpNegro}
\setbeamercolor{date}{fg=utpNegro}
\setbeamercolor{block title}{fg=white,bg=utpNegro}
\setbeamercolor{block body}{fg=utpNegro,bg=utpGrisClaro}
\setbeamercolor{alertblock title}{fg=white,bg=utpRojo}
\setbeamercolor{alertblock body}{fg=utpNegro,bg=utpRojo!7}
\setbeamercolor{exampleblock title}{fg=white,bg=utpGris}
\setbeamercolor{exampleblock body}{fg=utpNegro,bg=utpGris!8}
\setbeamercolor{terminal}{fg=white,bg=utpNegro}

\setbeamerfont{title}{series=\bfseries,size=\fontsize{19.5}{22.5}\selectfont}
\setbeamerfont{subtitle}{size=\fontsize{11.5}{13.5}\selectfont}
\setbeamerfont{date}{series=\bfseries,size=\fontsize{11.5}{13.5}\selectfont}
\setbeamerfont{author}{size=\fontsize{10.5}{12.5}\selectfont}
\setbeamerfont{frametitle}{series=\bfseries,size=\large}
\setbeamerfont{block title}{series=\bfseries,size=\small}
\setbeamerfont{block body}{size=\normalsize}

\setbeamertemplate{navigation symbols}{}
\setbeamertemplate{headline}{}
\setbeamertemplate{blocks}[default]
\setbeamertemplate{itemize item}{\color{utpRojo}\small\raise.15ex\hbox{$\blacktriangleright$}}
\setbeamertemplate{itemize subitem}{\color{utpGris}\scriptsize$\bullet$}
\setbeamertemplate{enumerate item}{\color{utpRojo}\insertenumlabel.}
\setbeamertemplate{background canvas}{\color{white}\rule{\paperwidth}{\paperheight}}

% Título de diapositiva: negrita negra a la izquierda, logo UTP a la derecha
\setbeamertemplate{frametitle}{%
  \begin{tikzpicture}[remember picture,overlay]
    \node[anchor=north east,inner sep=0pt]
      at ([xshift=-7mm,yshift=-5.5mm]current page.north east)
      {\includegraphics[width=17mm]{utp-logo.png}};
  \end{tikzpicture}%
  \vspace{1.2mm}%
  \begin{minipage}[t]{.80\textwidth}\raggedright\insertframetitle\end{minipage}\par
  \vspace{1.1mm}{\color{utpRojo}\rule{\textwidth}{1.1pt}}
}

% Solo numeración discreta; la plantilla UTP no lleva pie de página
\setbeamertemplate{footline}{%
  \hfill{\color{utpGrisMedio}\scriptsize\insertframenumber}\hspace{6mm}\vspace{2.5mm}%
}

% Portada UTP
\setbeamertemplate{title page}{%
  \begin{tikzpicture}[remember picture,overlay]
    \node[anchor=center,inner sep=0pt] at (current page.center) {%
      \includegraphics[width=\paperwidth,height=\paperheight]{utp-cover-bg.jpeg}%
    };
    \node[anchor=north,align=center,text=utpNegro,inner sep=0pt,
      text width=145mm,font=\usebeamerfont{title}]
      at ([yshift=-25.2mm]current page.north) {\inserttitle};
    \node[anchor=north,align=center,text=utpNegro,inner sep=0pt,
      text width=142mm,font=\usebeamerfont{subtitle}]
      at ([yshift=-38.2mm]current page.north) {\insertsubtitle};
    \node[anchor=north,align=center,text=utpNegro,inner sep=0pt,
      font=\usebeamerfont{date}]
      at ([yshift=-50.4mm]current page.north) {\insertdate};
    \node[anchor=north,align=center,text=utpNegro,inner sep=0pt,
      font=\usebeamerfont{author}]
      at ([yshift=-58.4mm]current page.north) {\insertauthor};
  \end{tikzpicture}%
}

% ---------- Comandos de apoyo didáctico ----------
\newcommand{\concepto}[2]{%
  \begin{block}{#1}\large #2\end{block}}
\newcommand{\pregunta}[1]{%
  \begin{alertblock}{Decide y explica}#1\end{alertblock}}
\newcommand{\analogia}[1]{%
  \begin{exampleblock}{En palabras simples}\small #1\end{exampleblock}}
\newcommand{\enlarealidad}[1]{%
  \begin{block}{Cómo se ve en la realidad}\small #1\end{block}}
\newcommand{\ojo}[1]{%
  \begin{alertblock}{Ojo con esto}\small #1\end{alertblock}}
\newcommand{\apoyo}{{\normalfont\footnotesize\color{utpGris}~· material de apoyo}}
\newcommand{\marca}[1]{\tikz[baseline=(m.base)]{%
  \node[rounded corners=1.5pt,fill=utpGrisClaro,draw=utpGrisMedio,inner xsep=3pt,inner ysep=1.5pt,
        font=\footnotesize] (m) {#1};}}
\newenvironment{terminal}{%
  \begin{beamercolorbox}[wd=\linewidth,sep=5pt,rounded=true]{terminal}%
  \ttfamily\small\ignorespaces}{%
  \end{beamercolorbox}}

\tikzset{
  caja/.style={draw=utpNegro, very thick, rounded corners=1.5pt, fill=white,
    align=center, minimum height=0.72cm, inner sep=5pt, font=\small},
  cajaazul/.style={caja, draw=utpGris, fill=white},
  cajaverde/.style={caja, draw=utpNegro, fill=utpGrisClaro},
  cajanaranja/.style={caja, draw=utpRojo!75, fill=utpRojo!7},
  flecha/.style={-{Stealth[length=2.3mm]}, very thick, draw=utpNegro},
  flechavuelta/.style={{Stealth[length=2.3mm]}-{Stealth[length=2.3mm]}, very thick, draw=utpNegro}
}

\graphicspath{{../../Sesion05/slides/assets/}}



\graphicspath{{assets/}{../../Sesion05/slides/assets/}}
\usepackage{colortbl}
\usepackage{listings}
\lstset{basicstyle=\ttfamily\footnotesize,backgroundcolor=\color{utpGrisClaro},
  frame=none,aboveskip=4pt,belowskip=4pt,showstringspaces=false,
  columns=fullflexible,keepspaces=true,breaklines=true,
  keywordstyle=\color{utpRojo}\bfseries,
  morekeywords={FROM,WORKDIR,COPY,CMD},
  commentstyle=\color{utpGris}\itshape,morecomment=[l]{\#},
  xleftmargin=5pt,xrightmargin=5pt,framesep=5pt}
% El código de Python de la práctica: solo comentarios y cadenas con color.
\lstdefinestyle{python}{basicstyle=\ttfamily\scriptsize,keywords={},
  stringstyle=\color{utpRojo},morestring=[b]",commentstyle=\color{utpGris}\itshape}
% Una terminal: fondo oscuro, sin colorear palabras.
\lstdefinestyle{terminal}{basicstyle=\ttfamily\scriptsize\color{white},backgroundcolor=\color{utpNegro},
  keywords={},morecomment=[l]{\#},commentstyle=\color{utpGrisMedio},escapeinside={(*}{*)},upquote=true}
\usepackage{upquote}
\usepackage{pifont}
\definecolor{okVerde}{HTML}{2E7D32}
\newcommand{\bien}{\textcolor{okVerde}{\ding{51}}}
\newcommand{\mal}{\textcolor{utpRojo}{\ding{55}}}
% Un número en un círculo, para señalar líneas de una terminal.
\newcommand{\nro}[1]{\tikz[baseline=(m.base)]\node[circle,fill=utpRojo,text=white,inner sep=1pt,font=\tiny\bfseries\sffamily](m){#1};}
% El estado de un contenedor, como etiqueta de color.
\newcommand{\estado}[2]{\tikz[baseline=(e.base)]\node[rounded corners=2pt,fill=#1,text=white,inner sep=2pt,font=\scriptsize\bfseries](e){#2};}

\title{Integración de Servicios Cloud GCP}
\subtitle{Docker desde cero: desplegar, imagen y contenedor}
\date{Semana 14 - Sesión 27}
\author{Rel Guzman}

\newcommand{\separador}[2]{%
\begin{frame}[plain]
\vfill
\begin{center}
{\fontsize{29}{34}\selectfont\bfseries #1}\par
\vspace{2mm}
{\small #2}\par
\vspace{4mm}
{\color{utpRojo}\rule{8.5cm}{1.2pt}}
\end{center}
\vfill
\end{frame}}

% Íconos simples para los diagramas.
\tikzset{
 caja/.style={draw=utpRojo,thick,rounded corners=4pt,align=center},
 vm/.style={draw=utpRojo,thick,rounded corners=3pt,fill=white,minimum width=1.6cm,minimum height=0.75cm,align=center,font=\scriptsize},
 vmmal/.style={vm,draw=utpGrisMedio,fill=utpGrisClaro,text=utpGris},
 zona/.style={draw=utpGrisMedio,dashed,rounded corners=4pt},
 flecha/.style={-latex,thick},
}
\newcommand{\cruz}[1]{\draw[utpRojo,line width=2pt] ($(#1)+(-0.25,-0.25)$) -- ($(#1)+(0.25,0.25)$) ($(#1)+(-0.25,0.25)$) -- ($(#1)+(0.25,-0.25)$);}

\begin{document}

{\setbeamertemplate{footline}{}\maketitle}

% 2
\begin{frame}{Revisión: lo que vimos en la Sesión 26}
\begin{columns}[T,onlytextwidth]
\column{.46\textwidth}
\begin{block}{Balanceador de carga}
\small Recibe a todos los clientes en una sola dirección y los reparte entre varios servidores iguales.
\end{block}
\begin{block}{CDN}
\small Guarda cerca del cliente una copia de los archivos que son iguales para todos.
\end{block}
\begin{block}{Contenedor}
\small Un paquete con tu programa y todo lo que necesita para funcionar.
\end{block}
\column{.50\textwidth}
\begin{alertblock}{Responde}
\small
1. Con 3 servidores y round robin, ¿a cuál va el cliente 5?\par\medskip
2. Si \texttt{web-2} no responde, ¿qué hace el balanceador?\par\medskip
3. ¿Por qué el carrito de Ana no va en la CDN?
\end{alertblock}
\end{columns}
\end{frame}

% 3
\begin{frame}{Revisión: la pregunta de hoy}
\begin{center}
\begin{tikzpicture}[font=\small]
  \node[caja,fill=white,text width=3.0cm,minimum height=1.2cm] (t) at (0,0) {\textbf{La tienda}\\\footnotesize código y bibliotecas};
  \node[caja,dashed,minimum width=1.6cm,minimum height=1.2cm,font=\Large\bfseries] (q) at (4.2,0) {?};
  \draw[flecha] (t) -- (q);
  \foreach \n/\y in {1/1.1,2/0,3/-1.1} {
    \node[vm] (w\n) at (8.0,\y) {\texttt{web-\n}};
    \draw[flecha,dashed] (q.east) -- (w\n.west);
  }
\end{tikzpicture}
\end{center}
\vspace{3mm}
\begin{columns}[T,onlytextwidth]
\column{.32\textwidth}
\begin{block}{1. Desplegar}
\small ¿Por qué copiar el código al servidor casi nunca alcanza?
\end{block}
\column{.32\textwidth}
\begin{block}{2. Imagen}
\small ¿Cómo se lleva la app junto con todo lo que necesita?
\end{block}
\column{.32\textwidth}
\begin{block}{3. Contenedor}
\small ¿Cómo se enciende, se detiene y se borra una copia?
\end{block}
\end{columns}
\end{frame}

% 4
\begin{frame}[c]
\begin{block}{Logro de la sesión}
Al finalizar, habrás encendido en tu computadora una \textbf{API} hecha con \textbf{Python} y \textbf{Flask}, y podrás explicar con tus palabras qué es \textbf{desplegar}, qué es una \textbf{dependencia}, en qué se diferencian una \textbf{imagen} y un \textbf{contenedor} y con qué comandos se enciende, se detiene y se borra un contenedor.
\end{block}
\end{frame}

% 5
\begin{frame}{Saberes previos}
\begin{itemize}\setlength\itemsep{14pt}
  \item En la Sesión 11, antes de encender el programa de Python, instalaste algo con \texttt{pip}. ¿Qué era y para qué servía?
  \item Cuando escribes \texttt{localhost:8080} en el navegador, ¿qué indica el 8080?
  \item Si cierras tu laptop, ¿alguien más puede usar un programa que tenías encendido en ella?
\end{itemize}
\end{frame}

% 6
\begin{frame}[c]{Utilidad de esta sesión}
\begin{center}
\begin{minipage}{.84\textwidth}
{\large Para que cualquiera use una aplicación, hay que \textbf{desplegarla}: llevarla a un servidor y dejarla funcionando. Hoy empezarás encendiendo una API para ver qué necesita para funcionar. Después verás cómo Docker la lleva junto con todo eso en una \textbf{imagen} y cómo se manejan sus \textbf{contenedores}.}
\end{minipage}
\end{center}
\end{frame}

% 7
\separador{Práctica}{La API de Café Nube, con Python y Flask en tu computadora}

% 8
\begin{frame}[fragile]{Práctica: la API de Café Nube}
\begin{block}{Caso de estudio}
\small \textbf{Café Nube}, una cafetería de barrio, quiere mostrar su carta en internet. Primero necesita una \textbf{API}: un programa que no muestra páginas, sino que \textbf{entrega datos} a otros programas.
\end{block}
\begin{columns}[T,onlytextwidth]
\column{.45\textwidth}
\begin{lstlisting}[style=python]
import os
from flask import Flask, jsonify
app = Flask(__name__)

MENU = [
    {"nombre": "Americano", "precio": 8},
    {"nombre": "Capuchino", "precio": 10},
    {"nombre": "Mocaccino", "precio": 12},
]

@app.get("/menu")
def menu():
    return jsonify(MENU)

puerto = int(os.environ.get("PUERTO", 8080))
app.run(port=puerto)
\end{lstlisting}
\column{.52\textwidth}
\renewcommand{\arraystretch}{1.2}
\scriptsize
\begin{tabular}{@{}>{\ttfamily}l>{\raggedright\arraybackslash}p{3.7cm}@{}}
\toprule
\textsf{\textbf{Línea}} & \textbf{En palabras}\\
\midrule
from flask import ... & Usa Flask, una biblioteca para hacer APIs\\
MENU = [...] & La carta: nombre y precio en soles\\
@app.get(\string"/menu\string") & Si alguien pide \texttt{/menu}\dots\\
return jsonify(MENU) & \dots le responde la carta en JSON\\
os.environ.get(...) & Lee la variable \texttt{PUERTO}; si no existe, usa 8080\\
app.run(port=puerto) & Enciende la API y espera pedidos\\
\bottomrule
\end{tabular}
\par\vspace{3mm}
{\scriptsize\textbf{JSON:} texto con datos que otro programa puede leer:\par \texttt{\{\string"nombre\string":\string"Americano\string",\string"precio\string":8\}}}
\end{columns}
\end{frame}

% 9
\begin{frame}[fragile]{Práctica: preparar y encender}
\begin{columns}[T,onlytextwidth]
\column{.56\textwidth}
{\small\bfseries Windows (PowerShell)}
\begin{lstlisting}
# 1. Instala Python
winget install --id Python.Python.3.12 -e
# 2. Cierra y abre PowerShell, y comprueba
python --version
# 3. Enciende la API (prueba 1)
python api.py
# 4. Instala lo de requirements.txt
python -m pip install -r requirements.txt
# 5. Enciéndela otra vez
python api.py
\end{lstlisting}
\column{.40\textwidth}
\renewcommand{\arraystretch}{1.25}
\scriptsize
\begin{tabular}{@{}>{\ttfamily}l>{\raggedright\arraybackslash}p{2.35cm}@{}}
\toprule
\textsf{\textbf{Comando}} & \textbf{Qué hace}\\
\midrule
winget install & Instala un programa en Windows desde PowerShell\\
--id ... -e & Elige el paquete con ese nombre exacto\\
python api.py & Ejecuta el programa \texttt{api.py}\\
python -m pip install & Instala las bibliotecas que pide \texttt{requirements.txt}\\
\bottomrule
\end{tabular}
\end{columns}
\par\vspace{3mm}
\begin{exampleblock}{Comprueba en el navegador: \texttt{http://localhost:8080/menu}}
\scriptsize\ttfamily [\{\string"nombre\string":\string"Americano\string",\string"precio\string":8\},\{\string"nombre\string":\string"Capuchino\string",\string"precio\string":10\},\par\ \{\string"nombre\string":\string"Mocaccino\string",\string"precio\string":12\}]
\end{exampleblock}
\par\vspace{2mm}
{\footnotesize\color{utpGris} Carpeta \texttt{Sesion27/practica\_cafe\_nube/}. La guía completa está en el \texttt{README.md}.}
\end{frame}

% 10
\begin{frame}{Práctica: seis pruebas}
\renewcommand{\arraystretch}{1.12}
\centering\small
\begin{tabular}{@{}c>{\raggedright\arraybackslash}p{5.9cm}>{\raggedright\arraybackslash}p{6.2cm}@{}}
\toprule
\textbf{\#} & \textbf{Qué hacer} & \textbf{Qué observar}\\
\midrule
1 & Paso 3: enciende la API \textbf{antes} de instalar \texttt{requirements.txt} & ¿Qué error sale? ¿Qué le falta a la API?\\
2 & Abre \texttt{localhost:8080/menu} o, en otra terminal, \texttt{curl.exe localhost:8080/menu} & ¿Qué responde? ¿Qué anota la terminal de la API?\\
3 & Abre \texttt{localhost:8080/postres} & ¿Qué código anota la terminal?\\
4 & En otra terminal, enciende la API otra vez & ¿Qué mensaje sale? ¿Por qué?\\
5 & En esa otra terminal: {\scriptsize\ttfamily \$env:PUERTO = \string"9090\string"} y luego \texttt{python api.py} & ¿Ahora sí enciende? Abre \texttt{localhost:9090/menu}\\
6 & Agrega un café a \texttt{MENU} y recarga \texttt{/menu} & ¿Aparece? ¿Y después de reiniciar la API?\\
\bottomrule
\end{tabular}
\par\vspace{4mm}
\begin{alertblock}{La idea}
\small Además de su código, la API necesita \textbf{Python}, \textbf{Flask} (lo que dice \texttt{requirements.txt}) y \textbf{un puerto libre}. Reto en el \texttt{README.md}: la ruta \texttt{/menu/<nombre>}.
\end{alertblock}
\end{frame}

% 11
\separador{Docker desde cero}{Desplegar, imagen y contenedor}

% 12
\begin{frame}{¿Qué es desplegar?}
\begin{center}
\begin{tikzpicture}[font=\small]
  \node[caja,fill=utpGrisClaro,text width=3.0cm,minimum height=1.5cm] (pc) at (0,0) {\textbf{Tu computadora}\\\footnotesize donde la escribes};
  \node[caja,fill=utpRojo!30,text width=3.0cm,minimum height=1.5cm] (srv) at (5.3,0) {\textbf{Un servidor}\\\footnotesize siempre encendido};
  \node[caja,draw=okVerde,fill=okVerde!12,text width=3.0cm,minimum height=1.5cm] (u) at (10.6,0) {\textbf{Cualquiera la usa}\\\footnotesize a cualquier hora};
  \draw[flecha] (pc) -- node[above,font=\footnotesize\bfseries,color=utpRojo]{desplegar} (srv);
  \draw[flecha] (srv) -- node[above,font=\footnotesize]{atiende} (u);
\end{tikzpicture}
\end{center}
\vspace{4mm}
\begin{columns}[T,onlytextwidth]
\column{.48\textwidth}
\begin{block}{Servidor}
\small Una computadora \textbf{siempre encendida y conectada} a la red para atender a otros. Lo que la hace servidor es su trabajo.
\end{block}
\column{.48\textwidth}
\begin{block}{Desplegar (\textit{deploy})}
\small Instalar y poner en marcha una app donde la van a usar, \textbf{para que quede funcionando}.
\end{block}
\end{columns}
\par\vspace{3mm}
\raggedright\small La API de Café Nube corre en tu laptop: si la cierras, \textbf{nadie más} puede pedir la carta.
\end{frame}

% 13
\begin{frame}{«En mi máquina funciona»: las dependencias}
\begin{center}
\begin{tikzpicture}[font=\small]
  \node[caja,draw=okVerde,fill=okVerde!10,text width=3.6cm,minimum height=2.3cm,align=left] (l) at (0,0)
    {\textbf{Laptop de Lucía}\\[2pt]\bien\ Python 3.12\\\bien\ Flask instalado\\\bien\ La app arranca};
  \node[caja,fill=utpRojo!8,text width=3.6cm,minimum height=2.3cm,align=left] (d) at (7.4,0)
    {\textbf{Laptop de Diego}\\[2pt]\mal\ Python 3.9\\\mal\ Sin Flask\\\mal\ La app falla al arrancar};
  \draw[flecha] (l) -- node[above,font=\footnotesize]{copia el código} node[below,font=\scriptsize\ttfamily]{api.py} (d);
\end{tikzpicture}
\end{center}
\vspace{4mm}
\begin{block}{Dependencia}
\small Todo lo que una app necesita para correr y \textbf{no es su propio código}: la versión del lenguaje, las bibliotecas, ciertos archivos del sistema.
\end{block}
\begin{alertblock}{Lo viste en la práctica}
\small \texttt{ModuleNotFoundError: No module named \textquotesingle flask\textquotesingle}: el código estaba completo, faltaba una dependencia.
\end{alertblock}
\end{frame}

% 14
\begin{frame}{Imagen y contenedor}
\begin{center}
\begin{tikzpicture}[font=\small]
  \node[caja,fill=utpGrisClaro,text width=2.7cm,minimum height=1.6cm] (app) at (0,0) {\textbf{Tu app}\\\footnotesize + Python, Flask y un pequeño Linux};
  \node[caja,fill=utpRojo!30,text width=2.4cm,minimum height=1.6cm] (img) at (4.6,0) {\textbf{Imagen}\\\footnotesize la plantilla};
  \draw[flecha] (app) -- node[above,font=\scriptsize]{empaquetar} (img);
  \foreach \n/\y in {1/1.0,2/0,3/-1.0} {
    \node[caja,draw=okVerde,fill=okVerde!12,minimum width=2.6cm,minimum height=0.7cm,font=\scriptsize] (k\n) at (9.6,\y) {\textcolor{okVerde}{$\blacktriangleright$} contenedor \n};
    \draw[flecha] (img.east) -- (k\n.west);
  }
  \node[font=\scriptsize] at (7.0,1.05) {\texttt{docker run}};
\end{tikzpicture}
\end{center}
\vspace{1mm}
\renewcommand{\arraystretch}{1.12}
\centering\small
\begin{tabular}{@{}l>{\raggedright\arraybackslash}p{4.2cm}>{\raggedright\arraybackslash}p{4.2cm}@{}}
\toprule
 & \textbf{Imagen} & \textbf{Contenedor}\\
\midrule
Qué es & Una plantilla: la app con sus dependencias & Una copia de la imagen \textbf{en marcha}\\
¿Corre? & Nunca & Mientras su programa trabaja\\
¿Cuántos? & Una & Todos los que quieras\\
\bottomrule
\end{tabular}
\par\vspace{2mm}
\begin{exampleblock}{Lo importante}
\small Si Lucía le pasa a Diego la \textbf{imagen}, Diego solo necesita \textbf{Docker}: Python y Flask ya vienen adentro.
\end{exampleblock}
\end{frame}

% 15
\begin{frame}[fragile]{Tu primer contenedor: \texttt{docker run hello-world}}
\begin{columns}[c,onlytextwidth]
\column{.60\textwidth}
\begin{lstlisting}[style=terminal]
$ docker run hello-world
Unable to find image 'hello-world:latest' locally (*\nro{1}*)
latest: Pulling from library/hello-world
Status: Downloaded newer image for hello-world:latest

Hello from Docker! (*\nro{2}*)
This message shows that your installation
appears to be working correctly.
\end{lstlisting}
\column{.36\textwidth}
\small
\nro{1}\ No tenía la imagen: \textbf{la descarga}. Solo la primera vez.\par\bigskip
\nro{2}\ Crea un contenedor, lo \textbf{arranca} y el programa escribe su saludo. Después termina.
\end{columns}
\par\vspace{5mm}
\begin{block}{\texttt{docker run} + una imagen}
\small Crea un contenedor de esa imagen y lo arranca. Si la imagen no está en tu computadora, primero la descarga.
\end{block}
\par\vspace{2mm}
{\footnotesize\color{utpGris} Antes, Docker Desktop debe estar abierto. Si responde \texttt{Cannot connect to the Docker daemon}, ábrelo y vuelve a intentar.}
\end{frame}

% 16
\begin{frame}{El ciclo de vida de un contenedor}
\begin{center}
\resizebox{\textwidth}{!}{%
\begin{tikzpicture}[font=\normalsize]
  \node[caja,fill=utpRojo!30,text width=2.5cm,minimum height=1.6cm] (img) at (0,0) {\textbf{Imagen} \texttt{nginx}\\\small la plantilla};
  \node[caja,draw=okVerde,fill=okVerde!12,text width=2.5cm,minimum height=1.6cm] (run) at (5.0,0) {\textbf{Corriendo}\\\small lo muestra \texttt{docker ps}};
  \node[caja,draw=utpGrisMedio,fill=utpGrisClaro,text width=2.5cm,minimum height=1.6cm] (stop) at (10.0,0) {\textbf{Detenido}\\\small solo lo muestra \texttt{docker ps -a}};
  \node[caja,draw=utpGrisMedio,dashed,text width=2.5cm,minimum height=1.6cm] (rm) at (15.0,0) {\textbf{Borrado}\\\small ya no existe; la imagen sigue};
  \draw[flecha] (img) -- node[above,font=\small\ttfamily]{docker run} (run);
  \draw[flecha] (run) -- node[above,font=\small\ttfamily]{docker stop} (stop);
  \draw[flecha] (stop) -- node[above,font=\small\ttfamily]{docker rm} (rm);
\end{tikzpicture}}
\end{center}
\vspace{4mm}
\begin{columns}[T,onlytextwidth]
\column{.48\textwidth}
\begin{block}{Segundo plano}
\small Una app de verdad no termina: queda esperando pedidos. Con \texttt{-d} corre \textbf{sin ocupar tu terminal}, y con \texttt{-{}-name web} le pones nombre.
\end{block}
\column{.48\textwidth}
\begin{block}{Logs}
\small Los mensajes que escribe el programa mientras corre. Se leen con \texttt{docker logs web} y son \textbf{lo primero que miras} cuando algo falla.
\end{block}
\end{columns}
\end{frame}

% 17
\begin{frame}{Los comandos de todos los días}
\renewcommand{\arraystretch}{1.35}
\centering\small
\begin{tabular}{@{}>{\ttfamily\footnotesize}l>{\raggedright\arraybackslash}p{5.4cm}c@{}}
\toprule
\textsf{\small\textbf{Comando}} & \textbf{Qué hace} & \textbf{\texttt{web} queda}\\
\midrule
docker run -d -{}-name web nginx & Lo crea y lo arranca en segundo plano & \estado{okVerde}{Corriendo}\\
docker ps & Lista los que corren: aparece \texttt{web} & \estado{okVerde}{Corriendo}\\
docker logs web & Muestra sus mensajes & \estado{okVerde}{Corriendo}\\
docker stop web & Lo detiene, sin borrarlo & \estado{utpGrisMedio}{Detenido}\\
docker ps -a & Lista también los detenidos: \texttt{Exited} & \estado{utpGrisMedio}{Detenido}\\
docker rm web & Lo borra; la imagen \texttt{nginx} sigue & \estado{utpRojo}{Borrado}\\
\bottomrule
\end{tabular}
\par\vspace{5mm}
\begin{alertblock}{Cuidado}
\small \texttt{docker rm} no borra un contenedor que está corriendo: primero \texttt{docker stop}.
\end{alertblock}
\end{frame}

% 18
\begin{frame}{Cierre: preguntas para comprobar}
\begin{block}{Responde con tus palabras}
\small
1. En la práctica, ¿por qué falló la API antes de instalar \texttt{requirements.txt}?\\\medskip
2. ¿Por qué no pudiste encender dos veces la API en el puerto 8080? ¿Cómo lo resolviste sin cambiar el código?\\\medskip
3. Lucía copia la API a un servidor, pero no la enciende. ¿Ya la desplegó?\\\medskip
4. Lucía le pasa a Diego la \textbf{imagen} de la API. ¿Qué necesita Diego para correrla?\\\medskip
5. Detuviste \texttt{web} con \texttt{docker stop}. ¿Sigue existiendo? ¿Con qué comando lo ves?
\end{block}
\begin{alertblock}{Idea final}
\small Una app no es solo su código: tu API necesita Python, Flask y un puerto libre. Una \textbf{imagen} lleva la app con todo eso, y cada \textbf{contenedor} es una copia en marcha que se enciende, se detiene y se borra.
\end{alertblock}
\end{frame}

% 19
\begin{frame}{Referencias}
\scriptsize
\begin{itemize}\setlength\itemsep{4pt}
  \item Práctica: \texttt{github.com/rgap/rgap-Integracion-de-Servicios-Cloud-Gcp-Agosto2026}, carpeta \texttt{Sesion27/practica\_cafe\_nube/}.
  \item Python: \texttt{python.org}. Flask: \texttt{flask.palletsprojects.com}. winget: \texttt{learn.microsoft.com/windows/package-manager}.
  \item Docker y Docker Desktop: \texttt{docs.docker.com}.
  \item Material previo del curso: Sesión 11 (Python, \texttt{pip} y variables de entorno) y Sesión 26 (imagen y contenedor).
  \item Las respuestas de las preguntas están en \texttt{../respuestas.md}.
\end{itemize}
\end{frame}

\end{document}
